\section{The operator selecting a good cut}

% BEGIN P84C1
\paragraph{P84C1. The polynomial labels of the finite central operator.}
At a generator prime $g_k$ and place $v$, put $A_{kv}=\mathbb F_{g_k}^r$.
The source indexes its local series by
\[
 P(a)=X^r+\sum_{j=0}^{r-1}a_j X^j,\qquad a\in A_{kv}.
\]
Each monic polynomial of degree $r$ occurs once, since its lower
$r$ coefficients recover $a$. This corrects LN1's description
of the indices as ordered root vectors. For instance $X^2+2$ over
$\mathbb F_5$ occurs at $(2,0)$ although it has no root in that
field. Even for split polynomials a root enumeration would introduce
multiplicities absent from this sum.

Let $E_{kv}(a)=e_{\rm Ser,k}(P(a))$ in the complex group algebra
of $G_{kv}=\mathrm{GL}_r(\mathbb F_{g_k})$. These are central
idempotents. The local finite type record sums its series over
their supported monic polynomials to $1$ and gives zero on the
other polynomials. By the bijection above,
$\sum_{a\in A_{kv}}E_{kv}(a)=1$. Embed each factor in the group
algebra of $G=\prod_{k,v}G_{kv}$ by extension at that place.
For a tuple $a$ write $E(a)$ for the product of the embedded
factors. Its coefficient at $g$ is the product of the local
coefficients at the coordinates of $g$. Thus the definition gives
\[
               e_0=\sum_{a:\,\mathrm{goodTuple}(P(a))} E(a),
                \qquad \sum_{a}E(a)=1.
\]
The second equality follows by finite distributivity. Each $E(a)$
is central and idempotent. The trace calculation of LN1 lifts
$|G|E(a)$ through the specified coefficient-ring embedding, using
the primitive $|G|$th root in the given depth-zero record. Summing
over the indicated tuples gives the central integral element $z_i$
with image $|G|e_0$.
% END P84C1

% BEGIN P84C2
\paragraph{P84C2. The selector acts on every type vector above a theta point.}
Assume the four compatibility conditions $\mathrm{Tie4}(D)$,
and take the automorphic constituent, twist and same-point clauses
of the Galois preconditions. Write $Y_i=Z_{\rm good}(i)$.
For every $x\in Z_i$, every automorphic representative $\Pi$
above $x$, and every vector $F$ of its specified type space,
\[
              e_0F=\mathbf 1_{Y_i}(x) F .
\]
First take $F$ in the parameter space of $\sigma$. At each $kv$
the $\mathrm{SerLoc}$ condition says that a local series acting
nontrivially has polynomial equal to $\mathrm{labPar}_{kv}(\sigma)$.
Its lower coefficients form the unique possible local vector $a_{kv}$.
If a tuple differs at any one place, that factor kills $F$, hence so
does its product. Summing all products gives $F$ by P84C1. Thus
$e_0F$ equals $F$ or zero according to the goodness of the
tuple of parameter labels.

It remains to put every $\Pi$ in the indicated parameter space.
At $x\in Z_i$ there is a representative with a base-change list
containing the distinguished twist of the character block
$\mu$ of length $t$. The same-point clause of
$\mathrm{provisionFortyOne}$ identifies the lists of two
representatives of $x$ up to permutation, and $\mathrm{adm1}$
supplies a list for $\Pi$. Move one occurrence of this block
to the end, preserving all multiplicities. Twist the remaining
list by $\eta_T$ and call the result $\Sigma$. Twisting
$\Sigma\mathbin{+\!+}[\mu]$ by $\eta_T^{-1}$ recovers the
permuted list: the twists are inverse on each positive-rank
constituent by $\mathrm{tie0}$ and $\mathrm{adm5}$. If
$x\in Y_i$, start with a good list $\Sigma$ witnessing that
membership. The same permutation argument transfers its
base-change identity to $\Pi$.

For either choice, parameter surjectivity in the
$\mathrm{Pend}$ witness of $\mathrm{Tie4}$ gives a parameter
$\sigma$ with these local values. The defining sum of subspaces
for $\mathrm{MtPar}(\sigma)$ contains the type of $\Pi$,
since $\Pi$ and $\Sigma$ satisfy its two incidence conditions.
The local polynomial dictionary identifies goodness of $\sigma$
with goodness of $\Sigma$ place by place. At $x\in Y_i$ our
chosen list is good. If $x\in Z_i\setminus Y_i$, a good
$\Sigma$ would itself witness $x\in Y_i$. The first paragraph
now proves the assertion for this very $F$.
% END P84C2

% BEGIN P84C3
\paragraph{P84C3. The good cut and the integral commuting operator.}
For every $D,\mathrm{nch}$ satisfying the native
$\mathrm{GaloisPreconditions}(D,\mathrm{nch})$ and
$\mathrm{Tie4}(D)$ there is a good cut, and the uncut datum
satisfies $\mathrm{CommutingEndomorphismCondition}$. Both
conclusions use the same operators. Keep $i$ exactly when
$Y_i$ is nonempty and set
\[
 S_i=\{x:\ \exists\Pi\text{ above }x,\ \exists F\in
       \operatorname{type}(\Pi),\ e_0F\ne0\}.
\]
Each point has a representative with nonzero type space. P84C2
gives $S_i\cap Z_i=Y_i$. With $N_G=|G|$, let
$\alpha_i=z_{O,i}(z_i)$ act on the integral Hecke module, for
$z_i$ of P84C1. Its action commutes with Hecke by the same
$\mathrm{Pend}$ witness. At $x\in Y_i$ its complex scalar on
all type vectors is $N_G$. At $x\notin S_i$ the scalar is zero,
including for $x$ outside $Z_i$, by definition of the support.
The $\mathrm{EigSpan}$ condition passes each scalar to the entire
rational eigenspace at $x$. This proves the operator property
of the cut and gives the commuting endomorphism at every tested
uncut index, since $Z_i\setminus Y_i$ misses $S_i$.

For the last goodness clause, fix $x\in S_i$ and a witnessing
$\Pi,F$. One good tuple has $E(a)F\ne0$. At each place its local
factor acts nontrivially, so one of the irreducible types in that
series occurs in $\Pi$. The goodness conditions give a squarefree
label $P(a)$ with nonzero constant term avoiding the prescribed
unit and its inverse at $j(i)$. For the root exclusions take
any auxiliary prime distinct from $g_k$ in the universal label
test. The $\mathrm{admDZ}$ field supplies zero monodromy and
the inertia polynomial with multiset from $P(a)(X-u)^t$ and its
conjugate-place prescription. A squarefree polynomial over the
perfect residue field is separable. Every repeated root of this
multiset is therefore $u$, and the corresponding unit at the
conjugate place also has order dividing $g_k-1$. The inertia
characters are trivial on wild inertia. Distinct nonzero labels
remain distinct on a tame generator: equality there gives
equality on the dense cyclic subgroup of the tame quotient and
then on inertia by continuity, where the characters distinguish
the labels. Multiplicativity thus transports the order condition
to every repeated inertia eigenvalue. Complete splitting and
the generator-prime size bound give residue norm $g_k$ and
multiplicity threshold two, yielding $\mathrm{Nc}$ as in LN2.

Finally take any other representative above $x$. Its base-change
list agrees up to permutation by the same-point clause. At these
split places the local parameter is given by this list, so it has
the same local condition. All five cut clauses follow, with
$\mathrm{Nc}$ on every representative. T83C3 gives $\mathrm{Tie4}$
on this same cut.
% END P84C3

% BEGIN P84C4
\paragraph{P84C4. The commuting condition on a good cut.}
In any good cut, at every retained index, the cut theta subset
equals the cut good subset. Indeed both memberships restrict
from the uncut point at that index. Goodness gives
$S_i\cap Z_i=Y_i$ and retains all of $Y_i$.
Take $\alpha=N_G\operatorname{id}_{M'}$ on the cut integral module.
This commutes with the restricted Hecke algebra and acts by the
required integer $N_G$ on every good eigenspace. The theta points
at which the commuting condition also asks for vanishing form
the empty set. This proves the commuting component of
$\mathrm{ExtraAlt}$ for every good cut, independently of the
norm-choice facts. It imposes no vanishing condition at the
retained points outside theta.
% END P84C4
